\documentclass[11pt,a4paper]{article}
\usepackage{fontspec}
\setmainfont{Noto Serif}
\setsansfont{Noto Sans}
\usepackage{babel}
\babelprovide[main,import]{russian}
\usepackage{amsmath,amssymb,mathtools,booktabs,geometry,fancyhdr,microtype,xcolor,hyperref}
\geometry{left=21mm,right=21mm,top=21mm,bottom=21mm}
\definecolor{ink}{HTML}{173D57}
\hypersetup{colorlinks=true,urlcolor=ink,linkcolor=ink,pdftitle={BGS-F. Критическое сжатие на двойных уровнях}}
\pagestyle{fancy}\fancyhf{}\lhead{\small BGS-F}\rhead{\small Критическое сжатие на двойных уровнях}\cfoot{\thepage}
\setlength{\headheight}{15pt}\setlength{\parindent}{1em}\setlength{\parskip}{2.5pt}
\setlength{\emergencystretch}{3em}\tolerance=3000
\newcommand{\Tr}{\operatorname{Tr}}
\newcommand{\EE}{\mathbb E}
\newcommand{\R}{\mathbb R}
\newcommand{\Z}{\mathbb Z}
\newcommand{\C}{\mathbb C}
\newcommand{\si}{\mathrm{Sine}_2}
\newcommand{\eps}{\varepsilon}
\newcommand{\sinc}{\operatorname{sinc}}
\newcommand{\sect}[1]{\section*{\sffamily\color{ink}#1}}
\newcommand{\subt}[1]{\par\medskip\noindent\textbf{#1.}\quad}
\newcommand{\T}{\mathbb T}
\newcommand{\dist}{\operatorname{dist}}
\newcommand{\Ad}{\operatorname{Ad}}
\newcommand{\Zp}{\mathbb Z_7}
\newcommand{\Qp}{\mathbb Q_7}
\newcommand{\Q}{\mathbb Q}
\newcommand{\HS}{\mathrm{HS}}
\begin{document}
\noindent{\small\sffamily\color{ink} МАТЕМАТИЧЕСКАЯ ФИЗИКА\hfill 22 НОЯБРЯ 2026}
\begin{center}{\LARGE\sffamily Критическое сжатие на двойных уровнях}\\[4pt]
{\large Совместный закон диагональных сдвигов и межузловых связей}
\end{center}
Пусть \(N=7^k\), \(m=3N/7\) и \(U_N(t)=e^{itA_N}M_N\) для исходных \(A_N,M_N\). Физическая сила \(u_N=t/N\). Берём выбранные линейные узлы с параметром характера \(\ell\equiv3\bmod7\). Они имеют кратность \(2\) и составляют \(1/7\) линейной сетки. Фаза \(j\)-го узла обозначается \(\lambda_j\), а переход к следующему выбранному узлу меняет \(j\) на \(j+7\).

\subt{Совместный закон целого семейства блоков}
Выбираем \(j\) равномерно среди этих узлов. После согласованной замены фаз базиса блоки \(\sqrt N A_N\) между узлами \(j+7r\), \(r\in\Z\), сходятся совместно на любом фиксированном наборе узлов. Предел является парой вещественных симметрических массивов \(K^\pm\). Для каждого \(r\le s\)
\[
 K^\pm_{rs}=\frac{14}{3}\Re\sum_{i\in I(s-r)}
              \eta_i c_i^{(\pm)} Z_{rs,i},
 \qquad (\eta_1,\eta_2,\eta_3,\eta_4)=(1,1,1,\sqrt7).
\]
Здесь \(Z_{rs,i}\) независимы и равномерны на единичной окружности, \(c_i^{(+)}=c_i\), \(c_i^{(-)}=\bar c_i\), а
\[
 (c_1,c_2,c_3,c_4)=\frac1{4000}(200-20i,-100+10i,-100+10i,1800).
\]
Множество \(I(d)\) содержит \(1,2,3\) всегда и содержит \(4\) тогда и только тогда, когда \(d\equiv0,3,4\bmod7\). Обе компоненты одной неупорядоченной пары делят одни и те же равномерные фазы. Разные пары независимы в пределе.

\subt{Критическая сила}
При \(u_N=\kappa N^{-3/2}\), \(\kappa\in\R\) фиксировано, получен предел выбранного сжатия точного \(U_N\) с двусторонними весами \((1+|r|)^{-1}\) в норме Гильберта--Шмидта по распределению. До наложения весов предельный массив равен
\[
       \operatorname{diag}\bigl(d_0 r I_2\bigr)+\kappa K
                      +i\kappa^2\frac{66033}{320000}I,\qquad d_0=98\pi/3.
\]
Для соответствующего конечномерного сжатия эрмитова спектрального карандаша Кэли мнимое скалярное слагаемое отсутствует.

\subt{Настоящие щели при убывающей силе}
При \(0<|u_N|=o(N^{-3/2}/\log\log N)\) нормированные внутриузловые щели настоящего \(U_N\) сходятся совместно на любом фиксированном наборе выбранных узлов к независимым копиям \(\tfrac7{300}|2Y_1-Y_2-Y_3|\), где \(Y_i\) независимы с арксинусным законом. При критической силе доказан операторный предел сжатия. Локальный спектр полного оператора на этой границе и BGS-F при фиксированной силе \(10^{-80}\) остаются открытыми.
\newpage
\sect{1. Индексы узлов и выбор ветвей}
Пусть \(\rho^2=2,\ \rho\equiv3\bmod7\), \(\gamma=3+2\rho\), \(\mathcal L(x)=\tfrac16\log(x^6)\) на \(\Zp^\times\). На \(\Qp^\times\) продолжаем \(\mathcal L\) по правилу \(\mathcal L(7)=0\) и аддитивности относительно умножения. Пусть \(\omega(3)=\zeta_3\) задаёт характер на \(\mathbb F_7^\times\). Как в предыдущих частях, метки \((v_1,v_2)\) переходят в расщеплённые координаты
\[
       (a,b)=\bigl(v_1+\rho v_2,\ (v_1-\rho v_2)/(-2\rho)\bigr).
\]
С точностью до общей фазы \(M_N\) есть дилатация \(\psi(x)\mapsto\psi(\gamma^{-1}x)\). Для \(\ell\equiv3\bmod7\) и \(h\bmod3\) оба вектора \(u_\pm\) имеют вид \(m^{-1/2}\omega(x)^h e_N(\ell\mathcal L(x))\) на квадратной или неквадратной орбите и ноль вне неё.

Поскольку \(\omega(\gamma)=\zeta_3^2\), связь параметров с номером узла \(j\) такова
\[
      j\equiv h\bmod3,\qquad
      j\equiv -3g\ell\bmod7^{k-1},\qquad
                         g=\mathcal L(\gamma)/7\in\Zp^\times .
\]
Здесь \(\nu_7(\mathcal L(\gamma))=1\) следует из порядка \(\gamma\) по модулям 7 и 49. Положим \(\delta=-7/(3g)\); тогда \(\delta/7\equiv3\bmod7\) по точному вычислению по модулю \(49\). У \(j+7r\) параметры \(X_r=X+r\delta\) и \(h_r=h+r\), где \(X=\ell\), а соответствующие остатки берутся по \(7^{k-1}\) и \(3\). Вся выбранная часть сетки есть один класс по модулю \(7\) среди узлов. Поэтому её можно однозначно нумеровать \(r\) в симметричном интервале длины \(m/7\) с центром \(r=0\).

В переходе из \(j+7s\) в \(j+7r\) положим \(d=s-r\), \(D=d\delta\). Для метки \((a,b)\) критическая фаза в единицах \(e_N\) имеет вид
\[
 \Phi(x)=b(x+a/2)+X_s\mathcal L(x+a)-X_r\mathcal L(x).
\]
Стационарное уравнение \(b x(x+a)+D x-X_r a=0\) проверяется по модулю 7. Для примитивных меток \(D\equiv0\bmod7\), поэтому прежний точный перебор для \(\ell\equiv3\) сохраняет только \(v=\pm v_i\), \(v_1=(0,1),v_2=(0,8),v_3=(0,-6)\). Все переходы между разными знаками квадратности исчезают. У каждой из оставленных меток по одному невырожденному корню на каждой орбите. При \(v_4=(0,7)\) и \(-v_4\) снимаем множитель 7, получая \(b_0 x(x+a)+(D/7)x-X_r a_0=0\). Дискриминант по модулю 7 равен \(2d^2+4\). Он ненулевой и является квадратом ровно при \(d\equiv0,3,4\). Произведение корней \(3\bmod7\) не является квадратом, так что есть один корень в каждой орбите. Разная квадратность невозможна, поскольку \(a\equiv0\bmod7\). Это доказывает правило \(I(d)\).

В каждом непрошедшем классе линейная фаза при разбиении \(x=x_*+7^{\lceil k'/2\rceil}y\) даёт нулевую геометрическую сумму, где \(k'=k\) или \(k-1\) после снятия \(7\). В прошедших классах гауссов расчёт предыдущей части даёт модули \(\sqrt N\) и \(\sqrt{7N}\) соответственно, с фазой, определяемой стационарным значением и характером. С учётом знаменателя \(m\) и множителя \(\sqrt N\) один такой вклад имеет модуль \(\tfrac73\eta_i|c_i|\).
\newpage
\sect{2. Общая фаза пары и вещественный базис}
Для \(v_i=(0,v_{i,2})\), \(q_i=a_i^{\rm W}b_i^{\rm W}=\rho v_{i,2}^2/2\), где \(a_i^{\rm W},b_i^{\rm W}\) --- расщеплённые координаты, введём \(z=x/a_i^{\rm W}\). Стационарное уравнение есть
\[
           q_i z(z+1)+D z=X_r.
\]
Два корня \(z_+,z_-\) соответствуют квадратному \(x\) для меток \(v_i,-v_i\) при \(x=a_i^{\rm W}z_+\) и \(x=-a_i^{\rm W}z_-\). Произведение двух \(x\)-корней у положительной метки равно \(-2\rho X_r\equiv3\bmod7\). Их квадратные классы различны. Для отрицательной метки корни переходят под \(x\mapsto-x\), и \(-1\) меняет квадратный класс. Для этих двух фаз
\[
 \Phi_{\pm,i}=q_i(z_\pm+1/2)+X_s\mathcal L(z_\pm+1)-X_r\mathcal L(z_\pm)
                         + D\mathcal L(a_i^{\rm W}).
\]
По Виету \(z_++z_-=-1-D/q_i\), \(z_+z_-=-X_r/q_i\), \((z_++1)(z_-+1)=-X_s/q_i\). Поэтому
\begin{equation}
   \Phi_{+,i}+\Phi_{-,i}
      =G(X_s)-G(X_r)+D\bigl(\mathcal L(2\rho)-1\bigr),
             \qquad G(x)=x\mathcal L(x).
                                                               \tag{1}
\end{equation}
Формула сохраняется у \(v_4\), так как \(\mathcal L\) продолжен на \(\Qp^\times\).

У двух гауссовых множителей произведение равно единице. Для примитивной пары гессианы по модулю 7 имеют противоположные знаки после подстановки \(x_-\equiv x_++a_i^{\rm W}\). У непримитивной пары после снятия \(7\) отношение коэффициентов квадратичных фаз принадлежит классу \(-1\). В знаменателе при \(a\equiv0\) остаётся квадрат, а производные стационарного многочлена на двух корнях противоположны. Квадратичные суммы с противоположными коэффициентами сопряжены, а замена коэффициента на тот же квадратный класс ничего не меняет. Произведение конечных характеров равно \(\omega(4)^d\). Действительно, \(x_+x_-=(a_i^{\rm W}/b_i^{\rm W})X_r=2\rho X_r\), и произведение конечных координат после сдвигов равно \(2\rho X_s\). Оба дают \(4\bmod7\). Соотношение степеней \(h_s-h_r=d\bmod3\) завершает расчёт.

Пусть \(V_{rs,i}\) --- полная единичная фаза вклада положительной метки из квадратной орбиты, вместе с гауссовым и конечным множителями. Введём узловую фазу
\[
   \Upsilon_r=\omega(4)^r e_N\!\bigl(G(X_r)-G(X)+r\delta(\mathcal L(2\rho)-1)\bigr).
\]
Вклад отрицательной метки имеет фазу \(\Omega_{rs}\overline{V_{rs,i}}\), где \(\Omega_{rs}=\Upsilon_s/\Upsilon_r\) независимо от метки. Чётность \(x\mapsto-x\) меняет две орбиты местами, меняет знак метки и сохраняет \(\chi_{\ell,h}(-1)=1\). Следовательно, в двух диагональных каналах блока между узлами коэффициенты положительных меток равны \(c_i\) и \(\bar c_i\) соответственно.

У каждого узла выбираем одинаковый для обеих орбит множитель базиса \(g_r=\Upsilon_r^{-1/2}\) с произвольной фиксированной ветвью. После этой замены блок веществен, поскольку фазовый множитель матричного элемента \(\bar g_r g_s\) имеет квадрат \(\Omega_{rs}^{-1}\). Новая фаза положительной метки есть \(\bar g_r g_s V_{rs,i}\), отрицательная является ей сопряжённой. Согласованность этой замены сразу на всех рёбрах опирается на точное тождество (1).
\newpage
\sect{3. Различные точки ветвления}
Выберем любой фиксированный конечный набор индексов \(r\), включив для удобства \(0\). Для каждого \(r\le s\) и \(i\in I(s-r)\) рассматриваем критическое значение \(\Phi_{+,i}^{rs}(X)\). Для \(r\ne0\) также включаем фазу \(B_r(X)=G(X_r)-G(X)+r\delta(\mathcal L(2\rho)-1)\). Все функции аналитичны на \(X\in3+7\Zp\) благодаря простым стационарным корням после снятия \(7\), где это необходимо. Их первые производные кратны \(7\), поскольку производная критического значения равна \(\mathcal L((x+a)/x)\) и \(B_r'=\mathcal L(X_r)-\mathcal L(X)\).

Для каждой критической фазы зависимость второй производной содержит квадратный корень
\[
 Y_i^{rs}=2q_i z_++q_i+D,\qquad
 (Y_i^{rs})^2=(q_i+D)^2+4q_i X_r.
\]
Точка ветвления над переменной \(X\) имеет вид
\begin{equation}
 b_{i,rs}=-\frac{q_i}{4}-\frac{r+s}{2}\delta
                               -\frac{(s-r)^2}{4q_i}\delta^2.
                                                               \tag{2}
\end{equation}
Эти точки различны для разных троек \(i,r,s\) при \(r\le s\). Для доказательства применяем \(7\)-адический вариант теоремы Эрмита--Линдемана. Если \(0<|z|_7<7^{-1/6}\) и \(z\) алгебраично над \(\Q\), то \(\exp_7 z\) трансцендентно. Здесь \(z=6\mathcal L(\gamma)\) ненулевое, \(\nu_7(z)=1\), а \(\exp_7 z=\gamma^6\) алгебраично. Значит, \(\mathcal L(\gamma)\), а потому и \(\delta\), трансцендентны над \(\Q\). Равенство двух правых частей (2) требует совпадения коэффициентов при \(1,\delta,\delta^2\). Первое даёт \(i=i'\) из четырёх различных \(q_i\). Последние два определяют неупорядоченную пару \(r,s\). Употребляется только теорема Малера об одном экспоненциальном значении. \href{https://ems.press/content/book-chapter-files/27396?nt=1}{Ю. Нестеренко, Mahler and Transcendence, §\,2.4, теорема 4}.

\subt{Независимость нелинейных фаз}
Дифференцируя по \(X\), имеем \(\Phi_{+,i}'=\mathcal L((z_++1)/z_+)\), \(z_+'=1/Y_i^{rs}\). Формулы Виета из раздела 2 дают
\begin{equation}
 (\Phi_{+,i}^{rs})''
 =\frac12\left(\frac1{X_s}-\frac1{X_r}\right)
   -\frac{2q_i X_r+D(q_i+D)}{2X_rX_s Y_i^{rs}}.
                                                               \tag{3}
\end{equation}
Коэффициент при \(1/Y_i^{rs}\) как рациональная функция ненулевой. Формула верна и для \(r=s\), где первый член исчезает.

В поле \(\Qp(X)\) квадратные классы \((Y_i^{rs})^2=4q_i(X-b_{i,rs})\) независимы, поскольку произведение любого непустого набора имеет нечётный порядок в некоторой точке \(b_{i,rs}\). Поэтому в их общем квадратичном расширении есть автоморфизм, меняющий знак одного выбранного корня и сохраняющий остальные. Это следует индукцией по числу радикалов. Если новый радикал уже содержался бы в прежнем расширении, он был бы собственным вектором всех замен знаков и равнялся бы рациональному множителю одного монома в прежних корнях, что нарушает независимость квадратных классов. Применяя отдельные замены знаков к комбинации вторых производных (3), исключаем все её коэффициенты при критических фазах, если комбинация постоянна. Оставшиеся \(B_r''=1/(X+r\delta)-1/X\) независимы по простым полюсам \(X=-r\delta\) для \(r\ne0\). Итак, никакая ненулевая целочисленная комбинация выбранных фаз не может быть постоянной.
\newpage
\sect{4. Совместный переход к равномерным фазам}
Рассмотрим ненулевую целочисленную комбинацию \(H(X)\) фаз из раздела 3. Производная \(H'\) аналитична и не тождественно нулевая на классе \(3+7\Zp\). Она не может обращаться в нуль на целой подокрестности: алгебраическая комбинация \(H''\) имела бы там тождественный нуль, вопреки различию ветвлений и простых полюсов. Поэтому на компактном классе \(3+7\Zp\) нули \(H'\) изолированы и конечны. В частности, доля \(\mathbb P(\nu_7(H')\ge n)\) стремится к нулю при \(n\to\infty\).

Значение \(e_N(H(X))\) зависит только от \(X\bmod7^{k-1}\), так как \(H'\) кратна \(7\). Точное разложение при \(X=X_0+7^{\lceil k/2\rceil}y\) даёт
\[
  H(X)\equiv H(X_0)+7^{\lceil k/2\rceil}y H'(X_0)\bmod7^k .
\]
Для примитивных корней и для корней после снятия \(7\) приращение \(x(X)-x(X_0)\) имеет интегральный ряд по \(X-X_0\). Логарифм относительного приращения отличается от своего линейного члена на \(O_7(7^{2\lceil k/2\rceil})\). Это обосновывает сравнение и то, что \(H'\) стабилен по модулю \(7^{\lfloor k/2\rfloor}\) в каждой подокрестности. Усреднение по \(y\bmod7^{\lfloor k/2\rfloor}\) даёт ноль, если \(H'(X_0)\not\equiv0\bmod7^{\lfloor k/2\rfloor}\). Оставшаяся доля стремится к нулю. По критерию Вейля все критические экспоненты и узловые фазы \(\Upsilon_r\), \(r\ne0\), сходятся совместно к независимым равномерным фазам. Гауссовы и конечные множители зависят только от паритета глубины, постоянных остатков \(r,s,i,h\) и опорного класса \(X\bmod7\), поэтому не меняют вывод. Он справедлив отдельно для каждого \(h\) и для их смеси.

Из раздела 2 реальные блоки получаются посредством умножения \(V_{rs,i}\) на \(\bar g_r g_s\) и добавления сопряжённой фазы. В пределе эта коррекция условно по узловым фазам есть поворот каждой независимой равномерной переменной, поэтому сохраняет совместную независимость. При выборе главного аргумента разрыв ветви квадратного корня имеет нулевую меру для предельных узловых фаз. Так получается вещественный массив с первой страницы. Его компоненты в двух каналах одной пары используют один и тот же набор \(Z_{rs,i}\). Для \(s=r+1\) вклад \(i=4\) отсутствует. Нормировка даёт
\[
 \mathbb E(K_{r,r+1}^{+})^2=\mathbb E(K_{r,r+1}^{-})^2
       =\frac{4949}{120000},
 \qquad \operatorname{Corr}(K_{r,r+1}^{+},K_{r,r+1}^{-})=\frac{99}{101}.
\]
Это следует из \(\mathbb E(\Re(c Z)^2)=|c|^2/2\) и
\(\mathbb E\Re(c Z)\Re(\bar c Z)=\Re(c^2)/2\), с тремя указанными коэффициентами.

Для малого \(t\) из предыдущего закона расщепления истинная приведённая щель на одном выбранном узле отличается от разности собственных чисел \(\sqrt N E_j A_N E_j\) на \(o_{\mathbb P}(1)\), если \(0<|t|\sqrt N\log\log N\to0\). Для каждого фиксированного \(r\) сдвиг \(j\mapsto j+7r\) переставляет выбранный класс, поэтому конечное объединение исключительных множеств сохраняет долю \(o(1)\). Поскольку диагональные пары нашего массива независимы для разных \(r\), совместный предел истинных внутриузловых щелей есть произведение их арксинусных законов.
\newpage
\sect{5. Точная критическая компрессия и скалярная утечка}
Пусть \(t=\kappa/\sqrt N\) и \(E\) --- сумма \(E_{j+7r}\) по любому фиксированному набору \(r\). В выбранном и приведённом к вещественным связям базисе
\[
   \frac{N}{i\lambda_j} E\bigl(U_N(t)-\lambda_j I\bigr)E
     = D_N+\kappa(\sqrt N EA_NE)L_N
                                +\frac{i\kappa^2}{2}(EA_N^2E)L_N+o(1),
                                                               \tag{4}
\]
где \(L_N=\lambda_j^{-1}M_N|_E\) имеет узловые фазы \(e_m(7r)\), а
\(D_N=N(L_N-I)/i\). Остаток по операторной норме \(O(N|t|^3)\to0\), поскольку \(\sup_N\|A_N\|<\infty\) по конечным коэффициентам кика. При фиксированном \(r\), \(D_N\to d_0 r I_2\), \(d_0=2\pi\cdot7\cdot(7/3)\).

\subt{Квадратичная часть}
Для каждого фиксированного ненулевого целого вейлева индекса \(v\) длина орбиты \(\Ad M_N\) на соответствующих метках имеет порядок \(N\): в расщеплённых координатах хотя бы одна компонента имеет конечную \(7\)-адическую валюацию, и \(\gamma\) имеет точные порядки по степеням 7. Проекция вейлева оператора \(T_N(v)\) на любое собственное подпространство \(\Ad M_N\) имеет ненормированный квадрат нормы Гильберта--Шмидта не больше \(N\) делённого на длину этой орбиты. С одной стороны, это следует разложением регулярного цикла в вейлевом базисе. С другой стороны, эта величина равна \(\sum_a\|E_a T_N(v)E_{a+b}\|_{\HS}^2\) для соответствующей разности \(b\). Поэтому по равномерному \(j\) из выбранной седьмой части любые два узла \(j+7r,j+7s\) дают
\begin{equation}
        \mathbb E_j\|E_{j+7r} T_N(v) E_{j+7s}\|_{\HS}^2
                                      \le C_v/N,
                                                               \tag{5}
\end{equation}
равномерно по \(r,s\), включая циклический переход.

Точное произведение Вейля разлагает \(A_N^2\) в конечную сумму с постоянной меткой \(\sigma^2=\sum_v|c_v|^2=66033/160000\). Остальные коэффициенты равномерно ограничены и находятся в фиксированном конечном носителе. По (5) \(EA_N^2E=\sigma^2 I_E+O_{\mathbb P}(N^{-1/2})\) для фиксированного набора узлов. Поэтому (4) стремится к \(\operatorname{diag}(d_0 r I_2)+\kappa K+i\kappa^2\sigma^2 I/2\). Мнимое слагаемое возникает из квадратичного члена точного унитарного оператора после сжатия.

\subt{Эрмитов карандаш Кэли}
Пусть \(C_t=\tan(t A_N/2)\), \(z=\lambda_j e^{ix/N}\) и \(-z\notin\operatorname{spec}(M_N)\). Точное тождество
\[
  (I-iC_t)(U_N(t)-zI)(M_N+zI)^{-1}
        =(M_N-zI)(M_N+zI)^{-1}+iC_t
\]
даёт, после деления на \(i\), эрмитов карандаш для линейной фазы \(z\). Его сжатие на \(E\) после умножения на \(2N\) сходится к \(\operatorname{diag}((d_0r-x)I_2)+\kappa K\), равномерно по \(x\) на компактных интервалах. Действительно, на конечном наборе выбранных узлов диагональный тангенс не имеет полюсов при ограниченном \(x\) и достаточно большом \(N\), а \(2N C_t=N t A_N+O(N|t|^3)\). Сжатый карандаш можно определять по выбранным диагональным элементам и в тех точках, где полный карандаш имеет полюс на другом узле.
\newpage
\sect{6. Взвешенный предел на всей выбранной части}
Пусть теперь \(E\) охватывает все \(m/7\) выбранных узлов. Индексируем их симметричным набором целых \(r\) длины \(m/7\) с центром \(0\) и сопоставляем пару базисных векторов \(\C^2\) на каждом \(r\in\Z\). Сжатые операторы продолжаются нулём за этим набором. Обозначим \(W=\operatorname{diag}((1+|r|)^{-1}I_2)\). Диагональная замена фаз из раздела 2 коммутирует с \(W\) и \(L_N\).

Определим \(\mathcal T_{N,\kappa}\) как (4) после умножения с обеих сторон на \(W\), в согласованном базисе. Тогда
\begin{equation}
  \mathcal T_{N,\kappa}\ \Longrightarrow\
    W\bigl(\operatorname{diag}(d_0 r I_2)+\kappa K
                               +i\kappa^2\sigma^2 I/2\bigr)W
            \quad\text{в пространстве } S_2 .
                                                               \tag{6}
\end{equation}
Здесь \(S_2\) --- операторы Гильберта--Шмидта. Предельный массив внутри весов задаётся поэлементно. Слева и справа весы делают его оператором Гильберта--Шмидта, поскольку \(K_{rs}\) равномерно ограничены, а \(r/(1+|r|)^2\) имеет суммируемый квадрат.

Для доказательства равномерная оценка (5) для всех фиксированных меток \(A_N\) даёт
\(\mathbb E_j\|\sqrt N E_{j+7r}A_NE_{j+7s}\|_{\HS}^2\le C\)
по всем \(r,s\) диапазона. То же доказательство для несмещённой части \(A_N^2-\sigma^2 I\) даёт \(C/N\) без множителя \(\sqrt N\). Поэтому за пределами фиксированного \(|r|,|s|\le B\) ожидаемый квадрат взвешенной нормы линейного члена ограничен
\[
 C_\kappa
       \sum_{\max(|r|,|s|)>B}\frac1{(1+|r|)^2(1+|s|)^2}
                                      =O_\kappa((1+B)^{-1}).
\]
Множители \(L_N\) не меняют абсолютные величины. Квадратичный член после отделения скаляра имеет ожидаемый квадрат нормы \(O_\kappa(N^{-1})\) на всём диапазоне. Диагональная часть сходится в \(S_2\), так как
\(\big|N(e_m(7r)-1)\big|\le d_0|r|\) и для конечного \(r\) есть поточечная сходимость. Взвешенный остаток \(O(N|t|^3)\) стремится к нулю в \(S_2\) из \(\|W\|_{S_2}<\infty\). На каждом конечном диапазоне применяем совместный закон раздела 4. Отсечение общего хвоста и затем \(B\to\infty\) дают (6), например по ограниченным липшицевым тестам в \(S_2\).

\subt{Численная и арифметическая проверка}
Целочисленный расчёт сверяет дискриминанты для всех семи классов разностей, общую фазу (1) и сумму \(\sigma^2\). Для ближайшей пары выбранных узлов одиннадцать точек (2) уже различаются по модулю 343. Доказательство произвольных конечных наборов использует трансцендентность \(\delta\). Отдельный матричный расчёт проверяет четыре узла при \(N=343,2401\), включая переключение непримитивного вклада, вещественность в согласованном базисе и тождество Кэли. Оба расчёта входят в архив.

\subt{Граница спектрального вывода}
В формуле (6) используется выбранная часть линейных узлов и два веса. Для собственных фаз полного \(U_N\) при критической силе требуется учесть исключённые узлы и взаимодействие дальних уровней в нерегуляризованной спектральной задаче. Из (6) не следует сходимость её корней. Исходная гипотеза с фиксированным \(10^{-80}\) остаётся открытой.
\end{document}
